\section{Model and Notation}

Cells are indexed by \(\F_2^4\), and square values are converted to labels by
subtracting one.  Thus the values \(1,\ldots,16\) become labels
\(0,\ldots,15\).  For a normal order-four magic square, let
\[
  L:\F_2^4\to\F_2^4
\]
be the label map.  Let \(A\) be the affine interpolation of \(L\) on the five
base points \(\{0,e_1,e_2,e_3,e_4\}\), and define the error map
\[
  E=L+A,\qquad \delta=|\supp(E)|.
\]

For four distinct selected labels, the selected set is an affine plane exactly
when its xor is zero.  This selected-xor axis is the finite selector used to
separate the exact affine row, B2, and B3 on endpoint \(24\).

\begin{definition}[Square-mask pair and endpoint]
A square-mask pair is a normal order-four magic square together with a
one-incidence mask.  The mask selects one cell in each row, one cell in each
column, and one cell on each of the two main diagonals.  The endpoint is
\(34-t\), where \(t=\min_{c\in M} L(c)\) is the minimum selected label (the
minimum selected value minus one).  Equivalently, endpoint \(24\) is exactly
the slice whose minimum selected value is \(11\).  This paper works inside the
certified endpoint \(24\) subatlas, which contains \(236\) such records.  In
particular every endpoint-\(24\) selected four-set contains the value \(11\),
which is why the B3 signatures are four-subsets of \(\{11,\ldots,16\}\)
containing \(11\).
\end{definition}

\begin{definition}[Selected mask and selected labels]
For a square-mask pair, \(M\subset \F_2^4\) denotes the four selected cells.
The selected labels are the four elements \(L(M)\subset \F_2^4\).  The selected
set is affine if and only if
\[
  \bigoplus_{c\in M}L(c)=0.
\]
Equivalently, in value notation, the selected values form an affine four-set
after subtracting one from each value.
\end{definition}

\begin{definition}[Endpoint-24 record]
An endpoint-24 record is a certified square-mask pair whose endpoint is \(24\).
The development atlas contains \(236\) such records.
\end{definition}

\begin{definition}[Cell-value affineness and essential representatives]
A record is cell-value affine when its full label map \(L:\F_2^4\to\F_2^4\) is
affine.  The finite atlas is stored modulo the usual dihedral \(D_4\) symmetries
of the square; theorem statements that count quotient objects explicitly say
when the free \(\TK\) quotient is also being used.
\end{definition}

\begin{definition}[The \(\TK\) quotient]
The finite transport group \(T\) is the row/column relation group used by the
extension certificates, and \(K\) is its loop stabilizer.  The verified quotient
\(\TK\cong V_4\) acts freely on every certified square-mask pair and preserves
the endpoint, affine-defect class, and selected signature.  Hence each
\(\TK\)-orbit has four records.
\end{definition}

\begin{definition}[Signature and support-action profile]
For B3, a signature is one of the eight non-affine selected-label four-sets
containing the minimum selected value \(11\).  A support-action profile records
the defect \(\delta\), the error-direction space \(\Derr\), and the finite
defect-mask support type seen by the selected mask.  The B3 support-action
matrix has these eight signatures as rows and the eight realized profiles as
columns.
\end{definition}

\begin{definition}[Clause templates]
The public names NF1, NF2, and NF3 refer to the three B3 support-action clause
templates.  NF1 is the defect-\(4\) singleton xor-error case.  NF2 is the
defect-\(6\) non-xor pair-support case.  NF3 is the defect-\(6\) non-xor
triple-support case.  The templates are structural; the thirteen profile-level
instances are finite-certified.
\end{definition}

\begin{definition}[Affine defect]
The affine defect of a record is \(\delta=|\supp(E)|\), where \(E=L+A\).  The
associated error direction space \(\Derr\) is the span of the bilinear error
coefficients.  The certified dichotomy is
\[
  (\delta,\dim \Derr)\in\{(0,0),(4,1),(6,2)\}.
\]
\end{definition}

\begin{definition}[Boundary labels]
Inside endpoint \(24\), we use the following names.
\begin{center}
\begin{tabular}{llrr}
\toprule
Name & Selector & Records & \(\TK\)-orbits\\
\midrule
exact affine core & selected affine, \(\delta=0\) & \(144\) & \(36\)\\
B2 & selected affine, \(\delta=4\) & \(32\) & \(8\)\\
B3 & selected non-affine, \(\delta\in\{4,6\}\) & \(60\) & \(15\)\\
\bottomrule
\end{tabular}
\end{center}
\end{definition}

The classical context is the order-four magic-square structure of Ollerenshaw
and Bondi~\cite{ollerenshawBondi}, vector-space viewpoints on magic
squares~\cite{wardMagicVectorSpaces}, the \Durer tesseract
interpretation~\cite{sudberyTesseract}, and enumeration work on magic squares
and the order-four case~\cite{beckCohenCuomoGribelyukMagic,
beckVanHerick4x4Enumeration}.  A recent
computational classification of the 880 essentially different normal
fourth-order magic squares is recorded by Takemura~\cite{takemura880}.  The
operative finite proofs in this paper are the local certificates listed in
Section~\ref{sec:certificates}.